\section{Related Work}
\label{sec:related}

\textbf{Localization that needs execution.} Spectrum-based and
mutation-based fault localization rank code by how it behaves under
failing and passing tests~\cite{just2014defects4j}. Recent agent-based
variants keep that dependency, using coverage or a reproduction
test~\cite{rafi2024multiagent,yu2025orcaloca}. The evidence they need
does not exist when an issue is filed. A measurement on crash reports
makes the point concretely: only 2 of 60 bugs had a fault-triggering
test available at report time~\cite{rafi2025sbest}. Our setting
excludes execution evidence by construction.

\textbf{Information-retrieval bug localization.} The BugLocator
lineage ranks files by textual similarity to the report, combining a
length-normalised vector-space score with SimiScore, a vote from
similar past reports~\cite{zhou2012buglocator}; BRTracer adds stack-trace
boosting and file segmentation~\cite{wong2014brtracer}; later work adds
learned rankers~\cite{ye2014learning} and deep
variants~\cite{niu2025deepirbl}. Bench4BL established the reproduction
baseline for this family~\cite{lee2018bench4bl}, and BLAZE extends it
across languages with dynamic chunking~\cite{chakraborty2025blaze}. We
re-implement rVSM, BugLocator and a BRTracer-style method for all four
of our languages and place them in the same study
(\cref{tab:study}). Two of our results concern this family directly.
Its similar-report vote is the same mechanism as our incident memory,
which is why we measure what that vote is worth once report-time
admissibility is enforced. And file segmentation, which helps the
vector-space scorers it was designed for, measured worse than whole-file
BM25 on our data, so we do not use it (\cref{sec:results:dense}).

\textbf{Logs as localization evidence.} Pathidea is the closest prior
work: it scores files using the stack traces and log lines in the report
and reconstructs execution paths between log statements to reach files
that logged nothing~\cite{chen2022pathidea}. Its log score and its
finding that gains track logging density inform our trace signal. Two
differences matter. Its path reconstruction is built on Java-specific
machinery, and three of our four projects are C or C++. And it does not
use the project's resolved history, which is the signal our study finds
the model cannot replicate.

\textbf{LLM and agent localization.} Agentless showed that a simple
two-stage procedure over the repository tree is a strong
baseline~\cite{xia2025agentless}, and it is the design we adopt for the
model-alone arm of our study. Agents improve on it by navigating
structure: LocAgent over a code graph with an entity
index~\cite{chen2025locagent}, CoSIL over dynamically built call
graphs~\cite{jiang2025cosil}, GenLoc by seeding an LLM with retrieval
candidates and letting it explore~\cite{asad2026genloc}. \sys{} follows
GenLoc's seed-then-explore shape and LocAgent's entity resolution. What
it adds is the project's own resolved history as a tool behind a
temporal firewall, and a design discipline these systems do not report
on: we measure which instructed behaviours the model actually performs,
and make the ones the result depends on structural rather than
instructed (\cref{sec:results:ablation}). CoSIL's analysis of LLM
localizer failure modes, in particular malformed tool arguments and
hallucinated paths, matches what we saw and shaped our tool interface.

\textbf{Incident history in operations.} Retrieving similar past
incidents is established practice in cloud operations. RCACopilot
retrieves prior incidents by embedding similarity with time decay and
reports that 93.8\% of recurring incidents reappear within 20
days~\cite{chen2024rcacopilot}; its vote weighting informs ours. RCAgent
performs tool-augmented root-cause analysis~\cite{wang2024rcagent}, and
LLM-based recommendation of root causes and mitigations is deployed at
scale~\cite{ahmed2023rootcause}. These systems identify a faulty service
or produce a prose explanation. They stop above the level of source
files, which is the step this paper takes.

\textbf{Benchmarks, leakage and contamination.} Our evaluation protocol
is shaped by three findings. Benchmarks built from public repositories
leak solutions: a third of SWE-bench instances were found to contain the
fix in the issue text or comments~\cite{aleithan2024swebenchplus}, and
models reproduce file paths for repositories they were trained on far
better than for held-out ones~\cite{liang2025illusion}. Issue-commit
link recovery inflates when future links are visible to the
model~\cite{niu2026ablots}, which is the leak our firewall closes and
measures. And localization accuracy drops sharply on data collected
after a model's training cutoff~\cite{zhang2025mulocbench}, while
time-aware rather than release-based evaluation changes which method
appears best~\cite{ledel2021broccoli}. We therefore use time-based
splits, a four-rule firewall, and a direct contamination probe
(\cref{sec:results:validity}).

\textbf{Industrial evaluation.} An evaluation of bug localization on a
large robotics system found simple term-frequency methods competitive
with transformer models and reported the practical constraints that
dominate adoption~\cite{hall2026industrial}. Our cost reporting and our
recommendation to retain a CPU-only fallback follow the same concern.
